\begin{center}
\zihao{3}\heiti 从强子关联函数和不变质量分布的角度研究奇特强子态的内部结构
\end{center}


\vspace{-0.2cm}
\centerline{\normalsize{\songti 研究生姓名：胡卓然}\hspace{40pt} \normalsize{ \songti 导师姓名：
\hspace{0pt}梁伟红~~~教授}}
%\centerline{\normalsize{\songti 研究生姓名：$\square\square\square$}\hspace{40pt} \normalsize{ \songti 导师姓名：
%\hspace{0pt}$\square\square\square$~~~教授}}
\centerline{\normalsize{ \songti 学科：物理学\quad}\normalsize{ \songti 研究方向： 粒子物理\quad}\normalsize{ \songti 年级： 2023 级} }
\vspace{0.5cm}
\centerline{\zihao{-3}\textbf{\songti 摘要}}
\vspace{0.1cm}

奇特强子态的结构和性质一直是强子物理中的重要研究课题，
对于理解强相互作用具有重要意义。
随着实验上越来越多的奇特强子态候选者被发现，由于现有的手段很难对其性质做出明确判断，因此发展有效的方法来检验这些态的性质显得尤为重要。
近期随着费米测量技术的发展，ALICE 和 STAR 合作组给出了新的实验数据—关联函数，并进一步提取了散射长度和有效力程，为研究奇特强子态提供了新思路。

本文首先从散射长度和有效力程的角度出发，探讨了其在检验奇特强子态性质中的潜力。首先，我们在手征幺正理论框架下研究了共振态 \(N^*(1535)\)，计算了其极点位置、散射长度和有效力程以及关联函数。随后将计算得到的可观测量作为来自实验的赝数据，在三夸克态图像下进行拟合。结果表明，三夸克态图像下提取的可观测量与分子态的预期结果存在显著差异，即可观测量能够有效区分不同的强子态图像。

在此基础上进一步讨论了能否通过关联函数精确提取可观测量，
并分别研究了轻味奇特强子态和重味奇特强子态。
对于轻味奇特强子态，我们研究了位于 \(\bar{K}N\) 阈值附近的奇特强子态候选者 \(\Sigma^*(1430)\)。在手征幺正理论框架下，我们计算了相关反应道的关联函数。采用最小模型依赖方法和重采样技术，我们成功提取了散射长度和有效力程，并且其不确定性低于 $20\%$，验证了从关联函数中高精度地提取可观测量的可行性。
在轻味奇特强子态研究成功的基础上，我们进一步探讨了该方法在重味强子体系中的应用，具体分析了 \(BD\) 相互作用。尽管束缚态位于阈值以下，关联函数仅能提供有限的束缚态尾部信息，但我们仍然成功提取了束缚态的信息，其结合能的不确定性约为 $50\%$，并以 $6\%$ 的精度确定了该束缚态的分子态概率。此外，散射长度的提取精度较高，而有效力程的不确定性相对较大，但仍与手征幺正理论的计算结果相符。这表明，关联函数同样适用于重味奇特强子态的研究。
因此从关联函数提取可观测量具有普适性，可以用以检验奇特强子态的性质。

本文还讨论了从不变质量分布提取可观测量的可能性，包括虚态的尖峰结构及阈下束缚态的近阈的不变质量分布。
针对长期存在争议的 \(K^{+}\bar{K}^{0}\) 的散射长度，我们利用 \(\chi_{c1} \to \eta\pi^+\pi^-\) 衰变中 不变质量分布，成功提取了 \(K^+\bar{K}^0\) 的散射长度和有效力程，并且不确定性低于 $20\%$。
此外通过分析$\Lambda_b \to \Lambda_c^+ \bar{D}^0 K^-(D^- \bar{K}^0)$ 衰变中$\bar{D}\bar{K}$阈值以上50 MeV的不变质量分布，成功提取了$D^{*}_{s0}(2317)$的质量（不确定性 4 MeV）及其分子态概率（不确定性 $11\%$）。

综上所述，本文的研究充分证明了可观测量在分析奇特强子态性质中的重要作用。我们不仅能够从关联函数中提取可观测量，在某些情况下还可以从不变质量谱中精确地提取相关信息。本文的研究为未来的实验研究提供了理论基础，随着实验测量的不断进步，我们有望进一步揭示奇特强子态的结构和性质。

\par
\vspace{0.5cm}
\noindent{\heiti 关键词：}强相互作用；强子结构；奇特强子态；关联函数；不变质量分布
